Toen Agape ontwaakte
en geleidelijk opkeek,
zag ze hoe zorg haar raakte
en liefde op haar neerstreek.

De trouwe gouvernante
ontving haar met een glimlach.
"Heb je gedroomd, Agape?",
zei ze toen ze haar blik zag.

"Dat is prachtig, lieve meid.
Dromen brengen ons verder!"
Liebling voelde zich bevrijd
en beschermd door haar herder.

Zo werd ze steeds meer verlicht
tijdens haar jonge leven.
Zonden kwamen aan het licht
door de duistere zeven

vrouwen die Lieblings mama
lopende haar zwangerschap
overdroeg als het coda
van haar korte ouderschap.

Haar moeder had ze onthuld
en intern getransformeerd.
Hun zendingswerk werd vervuld
en door Agape geleerd

via haar moeder Emmy.
Zo vond ze haar Ikigai,
haar pure levensmissie,
haar spiritueel karwei,

al voor ze werd geboren.
In haar werd mammie Emmy
als een boeddha herboren
als weg naar mitigatie

van moeders aardse zorgen.
Dit werd Liebling’s levenspad.
Doch een roeping voor morgen.
Ze gaf wat zij al bezat

in haar nog prille leven.
Ze zou nog meer gaan geven
aan vreugde en gemoedsrust als ze de enige lust

die haar missie bezielde
zou kunnen bevredigen.
Dat wat haar ooit vernielde
als plunje kon ledigen;

de bron van haar verwardheid.
Haar nog open dilemma,
haar onvoltooide eenheid;
de scheiding van haar mama.

Agape verkende angst
en kende daarvoor geen angst.
Dat was een goede vangst,
want die huist het langst.

Ze begaf zich meer en meer
op straat in haar vader’s stad.
Ze tergde iedere keer.
Stak een lont in zijn kruidvat.

De landheer was niet gewend
aan authentieke ondeugd.
Hij was stiekem wel verheugd
met zijn ferme dissident.

Ze loog nooit over een daad
of de hangplek des onheils.
Hij was bang, maar nooit echt kwaad.
Hij gunde het haar dikwijls,

want hij zag haar begaafdheid.
Hij kende ook het gevaar
en wilde Liebling niet kwijt
aan een wrokkige barbaar.

Op straat werd Liebling bekend
om haar kalm en wijs gezicht.
Ze werd voor velen een licht,
waarnaar men zich had gewend.

Zij luisterde geduldig
en wachtte tot haar woorden
vanzelf en zorgvuldig
haar stilzwijgen verstoorden.

Haar raadgevingen boeiden
en susten elke klager.
De antwoorden die vloeiden,
pasten bij elke vrager.

Als een wijze haar bezocht,
voelde die zich zelfs verstaan.
Als Orakel opgezocht,
trok zij steeds meer mensen aan.

Wanneer men Agape vroeg;
"Van waar komt jouw woordenpracht?",
dan kreeg de controleploeg:
"Niet ik, maar 't is de kracht."
"Spreek over hoge dingen
slechts met een verheven ziel.
Ga niet voor het volk zingen!”
sprak een geleerde schlemiel,

Meister Eckhart citerend.
"Wie begrijpt echt het leven,
maakt van inzicht iets lerend,
als hij dit niet kan geven

aan elk verlangend wezen?
God vertoont zich overal.
Aldaar kan men hem lezen.
Wie dit niet ziet, komt ten val,

struikelend in zijn verhaal.
God staat buiten elk begrip.
Liefde vang je niet met taal.”
was Agape’s slimme tip.

Bovenal was Agape
zelf haar verlichte boodschap,
die radiaal uitstraalde
als een halo rond haar kap.

Een magnetisch liefdesveld,
dat mensen naar haar toe trok.
Dieren met alle geweld
liet proberen uit hun hok

te gaan om bij haar te zijn.
Stralen duwden door wolken
om via hun zonneschijn
haar gezicht te bevolken.

Het enige aandenken,
aan Jeanne en de jager,
wou de landheer pas schenken
aan de jonge weetgrager,

als ze genoeg gerijpt was.
Het was Jeanne’s toverboek.
Toen Agape dit boek las,
bleek het gauw gesneden koek,

omdat ze waardevrij keek.
Ofschoon het voor geleerden
louter hocus-pocus leek,
als die het bestudeerden.

Het grimoire stond vol met
spreuken en kruidenkennis.
Inktschetsen van een skelet
van een mens of hagedis

annoteerden de natuur,
die het heksenboek beschreef.
De kunst van magie was puur
en volgens Liebling niet scheef.

Het occulte was voor haar,
slechts iets wat verborgen is.
Niet iets duisters, kwaads of raar,
maar slechts verboden kennis

van natuurlijke krachten,
die niet van onder kwamen.
Ook geen codex der nachten
waar eunjers de macht namen.
